\section{End-to-end architecture}\label{sec:arch}
Figure~\ref{fig:arch} shows the deployed system. A macOS host (Apple M4, 16\,GB) runs
Lima~\cite{lima} with Apple's Virtualization framework. Lima boots one Ubuntu~24.04 arm64 VM
(4 vCPU, 6\,GB, kernel 6.8), which supplies the Linux features that 5G software needs (SCTP,
TUN, netns, tc). Inside the VM, three network namespaces model the three administrative
domains:

\begin{itemize}
\item \textbf{\code{ran}}: the UERANSIM gNB and all UEs. Each UE PDU session appears as a
  \code{uesimtun$N$} interface. Because UEs live in their own namespace, downlink packets
  addressed to a UE cannot be short-circuited by the core's kernel; they \emph{must} pass
  UPF\,$\rightarrow$\,GTP-U\,$\rightarrow$\,gNB.
\item \textbf{root (core)}: the Open5GS NFs on loopback SBI addresses (HTTP/2, port 7777,
  through the SCP), MongoDB, and the two UPFs. A veth pair (\code{n3-core}/\code{n3-ran},
  10.10.0.0/24) carries both N2 and N3.
\item \textbf{\code{dn}}: the external data network (192.168.100.2), reached over the N6 veth.
  It hosts the iperf3 servers, the URLLC echo server and the IoT collector. A NAT rule also
  gives the UE pools real Internet access.
\end{itemize}

\begin{figure}[H]\centering
\resizebox{\linewidth}{!}{%
\begin{tikzpicture}[
  font=\sffamily\scriptsize, >=Stealth,
  nf/.style={draw=black!60, rounded corners=2pt, fill=white, minimum width=13mm, minimum height=6mm, align=center},
  ue/.style={nf, minimum width=15mm},
  upf/.style={nf, minimum width=17mm, minimum height=8mm, line width=0.9pt},
  lbl/.style={font=\sffamily\tiny, fill=white, inner sep=1pt, text=ink2, align=center},
  dom/.style={draw=black!25, rounded corners=4pt, inner sep=6pt}]
% ---------------- RAN namespace
\node[ue, draw=embb] (u1) at (0,3.2) {embb1\\\tiny 10.45.0.2};
\node[ue, draw=embb] (u2) at (0,2.3) {embb2\\\tiny 10.45.0.3};
\node[ue, draw=urllc] (u3) at (0,1.4) {urllc1\\\tiny 10.46.0.2};
\node[ue, draw=miot] (u4) at (0,0.5) {miot1/2\\\tiny 10.47.0.2--3};
\node[ue] (u5) at (0,-0.4) {multi\\\tiny eMBB+URLLC};
\node[nf, minimum height=22mm, minimum width=14mm] (gnb) at (2.6,1.4) {gNB\\[2pt]\tiny UERANSIM\\\tiny nr-gnb\\\tiny 10.10.0.2\\\tiny 3 S-NSSAIs};
\foreach \u in {u1,u2,u3,u4,u5} \draw[<->, dashed, black!60] (\u.east) -- (gnb.west);
\node[lbl] at (1.3,-0.95) {Uu\,$\approx$\,RLS / UDP 4997};
\begin{scope}[on background layer]
\node[dom, fill=bgran, fit=(u1)(u5)(gnb), label={[font=\sffamily\scriptsize\bfseries]above:netns \texttt{ran}}] (ran) {};
\end{scope}
% ---------------- Core
\node[nf] (amf) at (5.6,3.4) {AMF\\\tiny N2 10.10.0.1};
\node[nf] (smf) at (8.0,3.4) {SMF\\\tiny 3 S-NSSAIs};
\node[nf] (ausf) at (5.0,5.0) {AUSF};
\node[nf] (udm) at (6.4,5.0) {UDM};
\node[nf] (udr) at (7.8,5.0) {UDR};
\node[nf] (pcf) at (9.2,5.0) {PCF};
\node[nf] (nssf) at (10.6,5.0) {NSSF};
\node[nf] (nrf) at (10.6,3.4) {NRF\\\tiny + SCP, BSF};
\node[nf, fill=black!5] (db) at (7.8,6.2) {MongoDB\\\tiny subscribers};
\draw[black!50] (udr) -- (db);
\draw[black!45, line width=1.4pt] (4.6,4.25) -- (11.4,4.25);
\node[lbl, anchor=west] at (11.45,4.25) {SBI (HTTP/2 :7777 via SCP)};
\foreach \n in {ausf,udm,udr,pcf,nssf} \draw[black!45] (\n.south) -- (\n.south |- 0,4.25);
\foreach \n in {amf,smf,nrf} \draw[black!45] (\n.north) -- (\n.north |- 0,4.25);
\node[upf, draw=embb] (upfa) at (6.4,1.2) {UPF-A\\\tiny shared\\\tiny GTP-U 10.10.0.1};
\node[upf, draw=urllc] (upfb) at (9.0,-0.2) {UPF-B\\\tiny dedicated\\\tiny GTP-U 10.10.0.3};
\draw[<->, black!70] (smf) -- node[lbl, pos=0.45] {N4 PFCP\\127.0.0.7} (upfa);
\draw[<->, urllc] (smf.south east) to[out=-60, in=90] node[lbl, pos=0.5] {N4 PFCP 127.0.0.17} (upfb.north);
% N1/N2/N3
\draw[<->, thick] (gnb.north) |- node[lbl, pos=0.75, above=1pt] {N2 NGAP/SCTP :38412\\(carries N1 NAS)} (amf.west);
\draw[->, thick, embb] ([yshift=2pt]gnb.east) -- node[lbl, pos=0.55, above] {N3 GTP-U :2152} (upfa.west);
\draw[->, thick, urllc] ([yshift=-14pt]gnb.east) |- node[lbl, pos=0.75, below] {N3 GTP-U :2152} (upfb.west);
\begin{scope}[on background layer]
\node[dom, fill=bgcore, fit=(amf)(db)(nssf)(upfa)(upfb)(nrf), inner xsep=10pt,
  label={[font=\sffamily\scriptsize\bfseries]above:root namespace: Open5GS 5GC v2.8.0}] (core) {};
\end{scope}
% ---------------- tun + DN
\node[nf, minimum width=11mm, font=\sffamily\tiny] (t1) at (11.1,1.8) {ogstun\\10.45/16};
\node[nf, minimum width=11mm, font=\sffamily\tiny] (t3) at (11.1,1.0) {ogstun3\\10.47/16};
\node[nf, minimum width=11mm, font=\sffamily\tiny] (t2) at (11.1,-0.2) {ogstun2\\10.46/16};
\draw[embb] (upfa.east) -- (t1.west); \draw[miot] (upfa.east) -- (t3.west); \draw[urllc] (upfb.east) -- (t2.west);
\node[lbl, align=center] at (11.1,2.6) {tc HTB+netem\\(DL per slice)};
\node[nf, minimum height=26mm, minimum width=19mm, align=center] (dn) at (14.4,0.9)
  {\textbf{Data network}\\192.168.100.2\\[2pt]\tiny iperf3 :5201--5204\\\tiny UDP echo :9000\\\tiny IoT collector :9100\\\tiny NAT $\rightarrow$ Internet};
\draw[<->, thick] (12.0,0.9) -- node[lbl, above] {N6} node[lbl, below, align=center] {n6-core\\tc UL/slice} (dn.west);
\foreach \tn in {t1,t2,t3} \draw[black!60] (\tn.east) -- (\tn.east -| 12.0,0) -- (12.0,0.9);
\begin{scope}[on background layer]
\node[dom, fill=bgdn, fit=(dn), label={[font=\sffamily\scriptsize\bfseries]above:netns \texttt{dn}}] {};
\end{scope}
\end{tikzpicture}}
\caption{Deployed end-to-end architecture with all interfaces labelled. The colour of each
user-plane path matches the slice colour used in all plots (\slice{eMBB} blue, \slice{URLLC}
orange, \slice{mIoT} green). Everything runs inside one Ubuntu VM on the Mac host.}
\label{fig:arch}
\end{figure}

\paragraph{Interfaces.} \textbf{Uu} is replaced by UERANSIM's \emph{Radio Link Simulation}
(RLS): RRC/NAS and user data travel between \code{nr-ue} and \code{nr-gnb} as UDP datagrams on
port 4997. That is the channel-simulation mode used in place of radio hardware. It has no PHY,
no MAC scheduling and no fading; \S\ref{sec:limits} covers what that implies.
\textbf{N1} (NAS-5GS) rides inside NGAP. \textbf{N2} is NGAP over SCTP (gNB 10.10.0.2
$\rightarrow$ AMF 10.10.0.1:38412). \textbf{N3} is GTP-U over UDP 2152, terminating at
10.10.0.1 (UPF-A) or 10.10.0.3 (UPF-B). The SMF chooses the UPF per PDU session, and the
choice reaches the gNB as the UL tunnel address in PDUSessionResourceSetupRequest.
\textbf{N4} is PFCP between the SMF (127.0.0.4) and each UPF. \textbf{N6} is the veth to the
\code{dn} namespace. The service-based interfaces (\textbf{N8, N10, N11, N12, N13, N15, N22},
Nnrf, Nbsf) are HTTP/2 on port 7777, using indirect communication through the SCP.

\paragraph{End-to-end path of one URLLC packet.} The \code{urllc\_probe} sends from 10.46.0.2
on \code{uesimtun2} in \code{ran}. \code{nr-ue} carries it over RLS to \code{nr-gnb}, which
encapsulates it in GTP-U (TEID from the PDU session) towards 10.10.0.3. UPF-B decapsulates and
writes it to \code{ogstun2}; the kernel then routes it over \code{n6-core} (tc class 1:20)
into \code{dn}, where the echo server replies. The reply follows the same path in reverse:
\code{ogstun2} egress (tc HTB/netem) $\rightarrow$ UPF-B $\rightarrow$ GTP-U $\rightarrow$
gNB $\rightarrow$ RLS $\rightarrow$ UE.
